% Paper-ready source for inclusion in the ACC manuscript.
% Requires: \usepackage{amsmath,amssymb}
% Bibliography: ACC2027_SOTA_REFERENCES.bib

\section{Related Work and Positioning}
\label{sec:related_work}

\subsection{Adaptive and Learned Augmentation for Aerial Vehicles}

Adaptive augmentation of a trusted flight controller is an established control
strategy.  In particular, L1Quad augments a geometric $SE(3)$ controller with an
$\mathcal{L}_{1}$ state predictor, piecewise-constant uncertainty adaptation, and
low-pass filtering, and provides transient and ultimate tracking-error bounds
across a broad set of physical quadrotor uncertainty tests
\cite{wu2025l1quad}.  Geometric neural adaptive control has likewise used
online neural weights to compensate wind-induced uncertainty with boundedness
results \cite{bisheban2021geometric}, while Neural-Fly combines an offline learned
representation with low-dimensional online adaptation and stability guarantees
for flight in strong winds \cite{oconnell2022neuralfly}.  DATT combines learned
feedforward--feedback tracking with an online $\mathcal{L}_{1}$ disturbance
estimate on a Crazyflie platform \cite{huang2023datt}.  These works establish
that adaptive augmentation, online uncertainty estimation, and learned flight
compensation are not by themselves new contributions.

Residual learning is also well established.  Residual reinforcement learning
adds a learned action to a conventional controller to retain useful prior
structure \cite{johannink2019residual}.  For quadrotor attitude tracking,
HPA--PPO uses
\begin{equation}
  u = u_{\mathrm{PID}} + \Delta u_{\mathrm{RL}},
\end{equation}
where a history--preview, attention-enhanced PPO policy produces the residual
\cite{li2027hpappo}.  Its results include numerical, software-in-the-loop, and
Pixhawk/Gazebo hardware-in-the-loop tests; the vehicle dynamics remain simulated
in the latter.  Our work therefore does not claim the first learned residual
over a classical attitude controller.  It instead studies a locally adapted,
low-dimensional neural-state residual whose learning variables and actuator
authority remain directly inspectable.

\subsection{Neuromorphic Control, NEF, and PES}

The Neural Engineering Framework (NEF) and Nengo support dynamical-system
realizations, probing, and local online learning in neural populations
\cite{bekolay2014nengo}.  The prescribed error sensitivity (PES) rule updates a
decoder using a locally available activity and an externally supplied error.
PES-based adaptive residual compensation has previously been demonstrated on a
robotic arm, including CPU, GPU, and Loihi realizations
\cite{dewolf2016adaptive,dewolf2020nengo}.  Most importantly for the present
paper, Olivares et al. implemented a Nengo spiking PID controller in which PES
supports online gain self-tuning, and validated it on a free-flying DJI Tello
\cite{olivares2026selftuning}.  Consequently, this paper does not claim the first
PES, Nengo, or online neuromorphic UAV controller.

Our question is narrower: can paired, continuous membrane--gate dynamics provide
a useful stateful basis for a bounded residual around a fixed quaternion
controller, and what part of the improvement is caused by those intrinsic
dynamics versus online decoder adaptation?  This distinction is important
because the present ablations show that fixed and conductance-adaptive
membrane--gate variants are numerically indistinguishable in the amplified-gate
regime, whereas enabling PES produces the large improvement.  Thus, the current
evidence supports the membrane--gate state as the decoder basis and PES as the
principal learning mechanism; it does not support attributing the improvement to
conductance adaptation alone.

\subsection{Quaternion Attitude Control and Actuator Constraints}

Geometric and quaternion attitude control are mature foundations
\cite{lee2010geometric,fresk2013quaternion}.  We use the shortest-path quaternion
error rather than claim a new attitude representation.  With
$q_e=q_d\otimes q^{-1}$, the vector error is
\begin{equation}
 e_q=\operatorname{sgn}(q_{e,0})\operatorname{vec}(q_e),
 \qquad
 u_0=\operatorname{sat}(-K_qe_q-K_\omega\omega).
 \label{eq:baseline}
\end{equation}
The neural command is strictly an additive correction,
\begin{equation}
 u=\operatorname{sat}\!\left(u_0+\beta(s\odot\ell)\right),
 \label{eq:bounded_residual}
\end{equation}
where $\ell$ is the per-axis residual limit and $\beta$ is the authority factor.
This construction preserves the classical command when $\beta=0$ and permits
explicit auditing of saturation and rotor feasibility.  Control allocation under
quadrotor actuator constraints is itself an established topic
\cite{monteiro2016allocation}; our contribution is to expose the residual through
that interface rather than treating the learned output as an unconstrained torque.

For normalized axis command $\bar u_j$, paired directional currents are
\begin{equation}
 I_j^+=I_b+I_g\max(\bar u_j,0),\qquad
 I_j^-=I_b+I_g\max(-\bar u_j,0).
 \label{eq:paired_currents}
\end{equation}
Each channel carries coupled membrane and gate states,
\begin{align}
 \tau_v\dot v_i & =-(v_i-v_{\mathrm{rest}})+I_i
       -c_i(v_i-v_{\mathrm{reset}}),\\
 \tau_r\dot r_i & =\sigma\!\left(\alpha[v_i-(v_{\mathrm{th}}-\gamma r_i)]\right)-r_i,\\
 c_i & =\left(g_0+\mathbb{1}_{\mathrm{adapt}}g_1|I_i|\right)r_i^2,\qquad
 a_i=\operatorname{clip}(G r_i,0,1).
 \label{eq:cmg}
\end{align}
The six directional activities are decoded as
\begin{equation}
 s=\operatorname{clip}(Wa,-1,1),
 \quad y^*=\operatorname{clip}(-k_e e_q-k_\omega^*\omega,-1,1),
 \label{eq:decoder_target}
\end{equation}
and the online rule is
\begin{equation}
 W_{k+1}=\Pi_{\mathcal W}\!\left[
 W_k+\eta(y_k^*-s_k)a_k^\top-\lambda W_k\right].
 \label{eq:pes_update}
\end{equation}
Equations~\eqref{eq:paired_currents}--\eqref{eq:pes_update} make the roles
unambiguous: $(v,r)$ generate the neural state, $a$ is the decoder regressor, PES
adapts $W$, and $\beta\ell$ bounds the residual authority.

\subsection{Position of the Present Work}

The contribution is the joint architecture and its controlled audit: (i) a fixed
quaternion-PD backbone; (ii) paired continuous membrane--gate state dynamics;
(iii) online PES adaptation restricted to a low-dimensional decoder; (iv) bounded
additive residual authority propagated through actuator allocation and saturation;
(v) comparisons with LIF, ALIF, Izhikevich, and rate alternatives; and (vi) matched
direct and Nengo implementations.  MuJoCo supplies numerical rigid-body,
constraint, contact, and actuator dynamics \cite{todorov2012mujoco}, but the
fidelity of each airframe still depends on the identified inertial, propulsion,
sensor, delay, and saturation parameters.  Likewise, execution in Nengo is an
implementation result and is not, without deployment measurements, a claim of
neuromorphic-hardware energy efficiency.

The evidence should therefore be reported in descending order of strength:
cross-mission MuJoCo results on X2 and the explicitly labeled enhanced-propulsion
CF2 design study; mechanism ablations separating conductance from PES adaptation;
neuron-family and direct-versus-Nengo equivalence tests; then physical 3-DOF Hover
results, including the negative full3 and reduced-yaw outcomes.  This ordering
supports a precise claim: the method is a transparent, reproducible neuromorphic
residual mechanism with strong simulated gains and limited but real hardware
evidence, not yet a proof-backed replacement for robust adaptive control or a
free-flight validation.

